\documentclass[instructions.tex]{subfiles}
\begin{document}
 
\chapter{Travailler au salut sans relâche}
 
Quoique Dieu commence notre salut et qu'il l'achève par sa grâce, néanmoins ce que nous devons faire de notre côté pour être sauvés, est tellement notre ouvrage, que personne ne peut le faire pour nous, personne ne peut le faire sans nous, personne ne peut le faire que nous. Le salut en ce sens est notre affaire personnelle; il faut donc y travailler nous-mêmes: c'est l'ouvrage de toute la vie. Ce n'est pas assez d'y travailler une fois, il faut y travailler toujours, et, autant que la faiblesse humaine le permet, y travailler sans relâche.

Donner quelques momens à son salut, et donner tout le reste aux plaisirs, aux affaires du monde; s'occuper à tirer sa famille de la poussière, et n'agir que faiblement pour se tirer soi-même de l'enfer, c'est s'abuser. Quand toutes les autres affaires ne réussiraient pas, ce ne serait rien en comparaison de la perte du salut. Les autres affaires peuvent se réparer; mais, quand on a manqué son salut, c'est pour toujours. Or, pour manquer son salut, il ne faut qu'un moment.

\og Veillez sur vous en tout temps\fg, dit le Sauveur; si nous dormons sur l'affaire de notre salut, tout est à craindre. Jonas dormait lorsqu'on délibérait de le jeter à la mer. Samson dormait lorsqu'il fut surpris par les Philistins. Les hommes dormaient lorsque l'ennemi sema la zizanie dans leur champ. Ces figures nous apprennent qu'il ne faut rien négliger, puisque tous les momens peuvent décider de notre salut.

Quand même vous auriez fait de grands progrès dans la sainteté, ne vous relâchez point. \og Que celui qui croit être ferme\fg, dit saint Paul, \og prenne garde de tomber\fg. Saint Pierre tomba, pour s'être exposé. Un regard fit pécher David. La femme de Loth, pour un moment de curiosité, fut punie de Dieu. Apprenons, par ces exemples, à ménager tous les momens. On ne peut être trop sur ses gardes, quand on court risque d'être perdu pour jamais, dit saint Grégoire\footnote{Nulla satis magna securitas ubi periclitatur aeternitas. S.Gre.}. Vivez donc de telle sorte, que le soin de votre salut l'emporte sur toutes les autres choses; n'échappez aucune occasion de l'assurer par la pratique des vertus; que dans tout ce que vous faites, dans tout ce que vous projetez, une intention droite et la vue de plaire à Dieu, déterminent vos actions et vos desseins, afin qu'à votre mort vous alliez receuillir dans l'éternité le fruit de vos bonnes oeuvres.

Après tout, puisque Dieu veut sincèrement votre salut, et qu'il y pense à tout moment, vous seriez bien ingrat et étrangement ennemi de vous-même, de manquer de retour pour Dieu et de charité pour vous, en négligeant la chose que Dieu désire le plus, et qui vous intéresse uniquement.


\section{Résolutions}

\primo Je condamnerai souvent ma lâcheté, et je me précautionnerai contre ce relâchement, en me rappelant cette pensée: Si je me sauve tout est gagné pour moi; comme aussi, si je me perds, tout est également perdu pour moi. 

\secundo Je regarderai comme étant perdu, et comme devant être retranché de ma conduite, tout ce qui ne pourra pas me rapporter à la gloire de Dieu et à mon salut. 

\prayer{Affermissez-moi, Ô mon Dieu, dans la résolution de travailler à mon salut sans relâche.}

\end{document}
